\section{Level B: evidence}\label{sec:levelb}

Level~B decides. Items that pass the gate, or that are appealed, are trialled on respondents
who answer them mixed with validated items. The verdict comes from the responses alone.

\subsection{Ability and item screening}

Respondent $i$'s ability is estimated as $\theta_i=(T_i-\bar T)/\mathrm{sd}(T)$, the
standardized total score over a set of anchor items already certified free of DIF (30 in the
repository simulations). For each trial item, a two-parameter logistic model
(2PL)~\cite{birnbaum1968,lord1980}
\begin{equation}\label{eq:2pl}
  P(X_{ij}=1\mid\theta_i)=\sig\bigl(a_j(\theta_i-b_j)\bigr),
\end{equation}
is fitted by logistic regression on the fixed $\theta$, giving a discrimination $a_j$ and a
difficulty $b_j$. A fit that does not converge, for example under complete separation, is
treated as undetermined. Stage~1 of the pilot keeps an item if its point-biserial correlation
with $\theta$ is at least $0.20$, its fit converged, and $a_j\ge0.6$. A negative point-biserial
reveals an inverted answer key, and the simulations catch one. Two gaps are acknowledged in the
repository and remain open \openq. First, $\theta$ is a proxy on the standardized-total metric,
while the thresholds come from a literature that uses an IRT metric~\cite{hambleton1991}.
Second, no guessing parameter is fitted, although multiple-choice items with four options are
exactly where guessing matters (design decision D25).

\subsection{DIF with observed groups: calibration only}

With an observed group $g_i\in\{-1,+1\}$, DIF is tested by the logistic
model~\cite{swaminathan1990}
$\logit P(X_{ij}=1)=\beta_0+\beta_1\theta_i+\beta_2g_i+\beta_3\theta_ig_i$, rejecting at
$|\beta_2|>0.40$, and by the Mantel--Haenszel common odds ratio with the ETS
classes~\cite{mantel1959,holland1988mh,zieky1993}. Both need a per-respondent group label.
\Cref{req:anon} forbids that label in production, so these tests are compiled only into closed
calibration pilots with declared attributes (design decision D20). Production uses the latent
variant below.

\subsection{DIF on latent classes}

The anonymity-compatible detector is a mixture IRT model~\cite{rost1990,cohen2005mixture}. Each
respondent belongs to an unobserved class $Z_i\in\{1,\dots,G\}$ with proportions $\pi_g$, and
responses are independent given $(\theta_i,Z_i)$:
\begin{equation}\label{eq:mixture}
  P(X_{ij}=1\mid\theta_i,Z_i=g)=\sig\bigl(a_{jg}(\theta_i-b_{jg})\bigr),\qquad
  \ell(\cdot)=\sum_i\log\sum_g\pi_g\prod_j P(x_{ij}\mid\theta_i,g).
\end{equation}
The implementation tries $G\in\{1,\dots,4\}$, each with a shared discrimination
($a_{jg}=a_j$, uniform DIF) and with a per-class one. Every candidate is fitted from seeded
starts with an analytic gradient, and the candidate with the lowest
BIC~\cite{schwarz1978} is kept. The search stops once an extra class no longer lowers the BIC.
The item statistic is $\mathrm{DIF}_j=\max_{g,h}|b_{jg}-b_{jh}|$ over classes holding at least 5\%
of the respondents. An item is flagged if the selected fit converged, has at least two classes,
and $\mathrm{DIF}_j>1.0$. The literature value for this gap is $0.5$. The repository reports
that $0.5$ flags every item of a batch with one biased item, keeps $1.0$ as provisional, and
defers the choice to a planned false-positive study. Batches must hold at least two items and
3{,}000 respondents. Mixture likelihood ratios are not $\chi^2$-distributed under the
null~\cite{self1987,mclachlan2000}, so the BIC comparison is a heuristic, not a calibrated test.

\subsection{Identification: why a single biased item is invisible}

The repository simulation found that one distorted item in a batch of eight is invisible and
two are detected. The following two propositions explain this.

\begin{proposition}[One class-dependent item]\label{prop:single}\proved{}
Let $\theta$ be observed without error and $Z$ independent of $\theta$, and suppose only item
$j$'s response function depends on $Z$. Then the likelihood factorizes as
$\prod_{k\neq j}P_k(x_k\mid\theta)\cdot\bigl[\sum_g\pi_gP_{jg}(x_j\mid\theta)\bigr]$, so the
class parameters are identified only through the marginal response curve
$m_j(\theta)=\sum_g\pi_g\,\sig(a_{jg}(\theta-b_{jg}))$. For a symmetric two-class shift
($\pi=\tfrac12$, $b_{j\pm}=b_j\pm\delta$, common $a$), with $x=a(\theta-b_j)$,
\begin{equation}\label{eq:single-expansion}
  m_j(\theta)=\sig(x)+\tfrac12(a\delta)^2\sig''(x)+O(\delta^4)
  =\sig\bigl((1-\tfrac14(a\delta)^2)\,x\bigr)
   +\tfrac12(a\delta)^2\,\sig'(x)\bigl(\tfrac{x}{2}-\tanh\tfrac{x}{2}\bigr)+O(\delta^4).
\end{equation}
\end{proposition}
\begin{proof}
With only item $j$ depending on $Z$, summing over $Z$ acts on item $j$'s factor alone. The
expansion averages $\sig(x\pm a\delta)$, so odd orders cancel. It then uses
$\sig''(x)=-\sig'(x)\tanh(x/2)$ and $\sig((1+\eta)x)=\sig(x)+\eta x\sig'(x)+O(\eta^2)$ with
$\eta=-\tfrac14(a\delta)^2$.
\end{proof}

In words, a single class-dependent item looks, to second order, like a 2PL item with a lower
discrimination. The residual $\tfrac{x}{2}-\tanh\tfrac{x}{2}=x^3/24+O(x^5)$ vanishes to third
order where the item is most informative. The same flattening is produced by error in
$\theta$, so the two are confounded in practice.

\begin{proposition}[Pairs of class-dependent items]\label{prop:pair}\proved{}
Under conditional independence given $(\theta,Z)$,
\begin{equation}\label{eq:pair-cov}
  \Cov(X_j,X_k\mid\theta)=\sum_g\pi_g\,\bigl(P_{jg}(\theta)-m_j(\theta)\bigr)\bigl(P_{kg}(\theta)-m_k(\theta)\bigr),
\end{equation}
which for two classes equals $\pi(1-\pi)\,\Delta_j(\theta)\Delta_k(\theta)$ with
$\Delta_j=P_{j2}-P_{j1}$. Any one-class model implies $\Cov(X_j,X_k\mid\theta)=0$. Hence two items
whose class-dependence overlaps in $\theta$ leave a signature that no one-class model absorbs.
\end{proposition}
\begin{proof}
$\E[X_jX_k\mid\theta]=\sum_g\pi_gP_{jg}P_{kg}$ and $\E[X_j\mid\theta]=m_j$. Expanding the
covariance gives \eqref{eq:pair-cov}. For two classes, substitute $m_j=(1-\pi)P_{j1}+\pi P_{j2}$.
\end{proof}

Both signals are of fourth order in $\delta$ for small $\delta$, but their constants differ
greatly. \Cref{tab:information} gives the expected likelihood-ratio contribution
($2N\times$ the per-respondent information) at $N=3000$. For one item, this is the minimal
expected Kullback--Leibler divergence between the mixture curve and any 2PL. For a pair, it is
the conditional mutual information $I(X_j;X_k\mid\theta)$.

\begin{table}[t]
\centering\small
\caption{Expected likelihood-ratio contribution at $N=3000$ of one biased item (the part no 2PL
absorbs) and of one biased pair (conditional mutual information), for $a=1.25$, $b=0$,
$\pi=\tfrac12$ and exact $\theta$. For comparison, the BIC penalty per extra parameter is
$\ln 3000\approx8$. Source: \texttt{scripts/levelB\_information.py}.}
\label{tab:information}
\pgfplotstabletypeset[
  col sep=comma,
  columns={delta,lr_single,lr_pair,ratio},
  columns/delta/.style={column name={shift $\delta$},fixed,fixed zerofill,precision=2},
  columns/lr_single/.style={column name={one biased item},fixed,fixed zerofill,precision=3},
  columns/lr_pair/.style={column name={one biased pair},fixed,fixed zerofill,precision=2},
  columns/ratio/.style={column name={ratio},fixed,precision=0},
  every head row/.style={before row=\toprule,after row=\midrule},
  every last row/.style={after row=\bottomrule},
]{data/levelB_information.csv}
\end{table}

\begin{finding}\label{find:information}\empirical{}
A biased pair carries 245--306 times the information of a single biased item. At $\delta=0.9$
and $N=3000$, a single item contributes $0.5$ to the expected likelihood ratio, far below any
penalty, while a pair contributes $130$. The repository simulation, which uses a proxy $\theta$,
measures about $100$ per biased pair. For small shifts, pairwise information scales as
$\approx0.048\,\delta^4$ per respondent (the constant falls to $0.033$ at $\delta=0.9$, where the
classes start to separate), so halving the shift requires about sixteen times the respondents for
the same power. The production detector agrees. With exact $\theta$ it selects no mixture in 8
null batches and detects nothing in 8 batches with one biased item ($N\in\{3000,6000\}$), and it
flags exactly the biased items in all 8 batches with two, with no false flag.
\end{finding}

\Cref{prop:single,prop:pair} turn \cref{req:batch} from an empirical rule into a structural one.
The latent classes are identified by \emph{covariation between items} given ability, so they need
at least two items that share a distortion. A campaign to tilt the bank leaves this signature. An
isolated distortion does not, and neither does a single biased item repeated across batches that
never contain another item biased the same way.

\subsection{A confound: error in the ability proxy}\label{sec:proxy}

The detector conditions on $\hat\theta$, a standardized anchor total, not on $\theta$.

\begin{proposition}[Proxy-induced dependence]\label{prop:proxy}\proved{}
Let $\hat\theta$ be computed from anchor responses that are independent of the trial items
given $\theta$, and let every trial item's response function be non-decreasing in $\theta$.
Then, with no DIF at all,
\begin{equation}
  \Cov(X_j,X_k\mid\hat\theta)=\Cov\bigl(P_j(\theta),P_k(\theta)\bigm|\hat\theta\bigr)\;\ge\;0,
\end{equation}
with strict inequality whenever $\Var(\theta\mid\hat\theta)>0$ and both functions are strictly
increasing.
\end{proposition}
\begin{proof}
Conditional independence gives $\E[X_jX_k\mid\hat\theta]=\E[P_j(\theta)P_k(\theta)\mid\hat\theta]$
and $\E[X_j\mid\hat\theta]=\E[P_j(\theta)\mid\hat\theta]$. Two non-decreasing functions of the
same random variable are non-negatively correlated (Chebyshev's association inequality), and
strictly so if both are strictly increasing and the variable is not degenerate.
\end{proof}

Local dependence therefore exists whenever the proxy is imperfect, with or without DIF. A
mixture whose classes shift all items in the same direction absorbs part of it, as a ``higher
than measured'' and a ``lower than measured'' class. The likelihood ratio then grows linearly
in $N$, while the BIC penalty grows like $\ln N$, so for large enough samples the BIC selects a
mixture on data that contain none. The estimated gap of \emph{every} item then equals the size of
the common shift. This is the observed-score version of a problem known in DIF research:
matching on an unreliable total inflates DIF, which is why SIBTEST corrects the matching
variable by regression~\cite{shealy1993sibtest}.

\begin{table}[t]
\centering\small
\caption{Likelihood ratio of the two-class uniform-DIF mixture against one class, with $\theta$
from 30 anchors (proxy) or exact, for batches of 8 items of which 0, 1 or 2 are biased
($\delta=0.9$). Mean of three datasets. Last row: BIC penalty $8\ln N$ used by the repository
simulation. Source: \texttt{scripts/levelB\_proxy.py}.}
\label{tab:proxy}
\begin{tabular}{@{}llrrrr@{}}
\toprule
$\theta$ & biased items & $N=1500$ & $3000$ & $6000$ & $12000$\\
\midrule
proxy & 0 & 53.7 & 91.2 & 169.0 & 281.9\\
proxy & 1 & 48.2 & 84.5 & 152.7 & 239.3\\
proxy & 2 & 67.1 & 219.7 & 433.4 & 568.8\\
\midrule
exact & 0 & 15.5 & 17.8 & 15.9 & 20.0\\
exact & 1 & 13.6 & 12.1 & 18.9 & 18.2\\
exact & 2 & 47.8 & 170.2 & 324.1 & 419.1\\
\midrule
\multicolumn{2}{@{}l}{penalty $8\ln N$} & 58.5 & 64.1 & 69.6 & 75.1\\
\bottomrule
\end{tabular}
\end{table}

\Cref{tab:proxy} shows the effect in a re-implementation of the repository simulation's model.
With exact $\theta$ the likelihood ratio stays flat when there is nothing to find, and also when
there is one biased item, as \cref{prop:single} predicts. With the proxy it grows linearly in
$N$ even with zero biased items, and it exceeds the penalty from $N=3000$ on. The mean
estimated gap of clean items is then $0.71$--$0.86$.

To measure the consequence for the production detector, we ran the Rust engine and the
protocol's verdict rule on null batches (no biased item, $N=6000$) whose $\theta$ comes from 10,
20, 30 or 60 anchors (\cref{tab:anchors}).

\begin{table}[t]
\centering\small
\caption{Production detector (Rust engine and verdict rule) on null batches: eight clean items,
$N=6000$, four datasets per row, $\theta$ from the stated number of anchors. A flag here is a
false positive. Source: \texttt{scripts/levelB\_detector.py}.}
\label{tab:anchors}
\begin{tabular}{@{}rcccc@{}}
\toprule
anchors & mixture selected & largest clean gap & batches with a false flag & items flagged per batch\\
\midrule
10 & 4/4 & 1.27--1.41 & 4/4 & 8 of 8\\
20 & 4/4 & 1.04--1.16 & 4/4 & 1--3\\
30 & 4/4 & 0.94--0.99 & 0/4 & 0\\
60 & 0/4 & --- & 0/4 & 0\\
\bottomrule
\end{tabular}
\end{table}

\begin{finding}\label{find:proxy}\empirical{}
On null batches the production detector selects a spurious mixture whenever the proxy is
imperfect enough. With 10 or 20 anchors it flags clean items in every batch. With 30 anchors,
the configuration of the repository simulations, it flags none, but the largest clean gap
reaches $0.94$--$0.99$ against a threshold of $1.0$. The provisional threshold is therefore
calibrated, implicitly, to one anchor configuration. With the same 30 anchors and one biased
item at $N=6000$, the detector flags the biased item in 2 of 4 batches (gaps $1.15$ and $1.03$).
It does so because the item rides on the artefactual class, not because its own shift is
identified; with exact $\theta$ it is never detected.
\end{finding}

Three remedies are available \openq.
\begin{enumerate}[label=(\alph*)]
  \item Put the anchors inside the likelihood with class-invariant parameters and integrate
    $\theta$ out, which gives a full mixture IRT model. The anchors then identify class-specific
    ability distributions, and a class-wide shift is no longer mistaken for DIF. This is the
    principled fix. Its cost is computational.
  \item Allow a class-specific ability shift and measure each item's departure from the batch's
    common class shift. In the two-class parametrization $b_{j\pm}=b_j\pm\delta_j$ of the
    repository simulation this is $2|\delta_j-\operatorname{median}_k\delta_k|$. In our re-implementation
    this ``differential gap'' separates biased items (at least $1.65$) from clean ones (at most
    $0.32$) at both sample sizes, with 30 anchors
    (\texttt{scripts/levelB\_differential.py}). A campaign that biases most of a batch in
    one direction would, however, be absorbed into the common shift, so this can only
    complement (a).
  \item At a minimum, make the threshold a documented function of anchor reliability and
    calibrate it by simulation, as the evaluation plan requires.
\end{enumerate}

\subsection{Sample sizes and the minimum network}

The pilot sizes of about $300$ (screen) and $1500$--$3000$ (DIF) are calibration targets taken
from the psychometric literature, not derived results. The repository labels them a
hypothesis \openq. \Cref{find:information} gives a first handle: for small shifts a biased pair
contributes an expected likelihood ratio of about $2N\times0.048\,\delta^4$ for $a\approx1.25$. Detection needs
this to exceed the penalty of the extra class, which is roughly $(K+1)\ln N$ for a batch of $K$
items. The respondent requirement sets a floor on the network's size. A batch needs a few
thousand \emph{distinct} respondents in its validation window, and anonymity is weak in a small
crowd. The repository estimates that the network cannot run as specified below about $2{,}000$
active participants.

\subsection{DIF is not bias}\label{sec:dif-not-bias}

DIF shows that, at equal measured ability, an item's difficulty depends on something else: the
item measures a secondary dimension on which the classes differ~\cite{shealy1993sibtest}.
Whether that secondary dimension is a nuisance is a judgment, not a statistical fact. For civic
knowledge the judgment matters. Suppose one camp is systematically misinformed about a
true fact. Then an item stating that fact shows DIF by construction, and Level~B rejects it for
the same reason Level~A did. The appeal to evidence therefore recovers a true but divisive item
only when the division lies in \emph{opinions about the item}, and not when it lies in
\emph{knowledge of the fact}. The design currently treats every politically aligned secondary
dimension as a nuisance. That is a normative choice, and we argue it should be made explicitly
\designchoice. The same mechanism detects distortions on axes nobody chose to look for. In the
repository simulation, an item neutral on the political axis shows DIF of $0.66$ on education.
Latent classes can surface such an axis, but only when a batch contains at least two items
distorted along it.
